\documentclass{article}
\usepackage[utf8]{inputenc}
\usepackage[T1]{fontenc}
\usepackage{amsmath}
\usepackage{graphicx}
\usepackage{tikz}
\usetikzlibrary{shapes.geometric,arrows,positioning,calc,er}
\usepackage{booktabs}
\usepackage{hyperref}
\usepackage{url}

% ---- Custom macros (deeply nested) ----
\newcommand{\Rel}[1]{\text{Rel}(#1)}
\newcommand{\Proj}[2]{\pi_{#1}\left(#2\right)}
\newcommand{\Select}[3]{\sigma_{#1}\left(#2,#3\right)}
\newcommand{\Join}[2]{{#1}\bowtie{#2}}
\newcommand{\semijoin}[2]{{#1}\ltimes{#2}}
\newcommand{\naturaljoin}[2]{{#1}\Join{#2}}

\graphicspath{{./figures/}}

\title{Database Systems: Entity-Relationship Modeling and Relational Algebra}

\author{Author Name}

\begin{document}

\maketitle

\begin{abstract}
We present a comprehensive study of database modeling techniques, covering
entity-relationship diagrams and relational algebra. Our analysis demonstrates
the application of these concepts to real-world scenarios.
\end{abstract}

\section{Entity-Relationship Modeling}

Entity-Relationship (ER) diagrams provide a visual representation of data
structures in database systems.

% ---- TikZ ER Diagram ----
\begin{figure}[ht]
  \centering
  \begin{tikzpicture}[
    node distance=2.5cm,
    entity/.style={rectangle,draw,fill=blue!10,minimum width=2.5cm},
    attr/.style={ellipse,draw,fill=gray!10,minimum width=1.8cm},
    rel/.style={diamond,draw,fill=green!10,minimum width=2cm,aspect=2}
  ]
    \node[entity] (student) {Student};
    \node[attr,above of=student] (sid) {\underline{ID}};
    \node[attr,left of=sid] (sname) {Name};
    \node[attr,right of=sid] (smajor) {Major};
    \draw (sid) -- (student);
    \draw (sname) -- (sid);
    \draw (smajor) -- (sid);

    \node[entity,right=3cm of student] (course) {Course};
    \node[attr,above of=course] (cid) {\underline{Code}};
    \node[attr,right of=cid] (ctitle) {Title};
    \draw (cid) -- (course);
    \draw (ctitle) -- (cid);

    \node[rel,right=1.5cm of student] (enroll) {Enroll};
    \node[attr,below of=enroll] (grade) {Grade};

    \draw (student) -- node[above] {N} (enroll);
    \draw (course)  -- node[above] {M} (enroll);
    \draw (enroll)  -- (grade);
  \end{tikzpicture}
  \caption{Entity-Relationship diagram for a university database.}
  \label{fig:er}
\end{figure}

\section{Relational Algebra}

Relational algebra provides theoretical foundations for query languages.

% ---- Relational algebra formulas (heavy custom macros) ----
\begin{align}
  &\Proj{\text{Name},\text{Grade}}(
    \Select{\text{Grade} > 3.5}{\sigma}{\Join{\sigma}{\text{Student}}}
  ) \label{eq:ra} \\
  &\semijoin{\Rel{\text{Enroll}}}{\Rel{\text{Course}}}
    \;=\;
    \Proj{\text{SID}}{\Rel{\text{Enroll}}}
    \;\cap\;
    \Proj{\text{SID}}{\Join{\Rel{\text{Enroll}}}{\Rel{\text{Course}}}} \label{eq:semijoin}
\end{align}

\section{Query Examples}

\begin{table}[ht]
\centering
\caption{Sample relational algebra queries}
\label{tab:queries}
\begin{tabular}{lp{6cm}}
  \toprule
  Query & Expression \\
  \midrule
  Q1 & $\Proj{\text{Name}}{\Select{\text{Major}=\text{CS}}{\text{Student}}}$ \\
  Q2 & $\Proj{\text{Title}}{\naturaljoin{\text{Enroll}}{\text{Course}}}$ \\
  Q3 & $\Proj{\text{Name}}{\semijoin{\text{Enroll}}{\text{Course}}}$ \\
  \bottomrule
\end{tabular}
\end{table}

\section{Conclusion}

ER modeling and relational algebra are fundamental tools for database design
and query formulation.

\bibliographystyle{plain}
\bibliography{refs}

\end{document}
